\section{Conclusion}
\label{sec:conclusion}

Preference is not tattooability on the pairs we audited.
On HPDv2 tattoo-subject pairs, pre-registered probes that penalise dense, shaded, thin-line images agree with human winners below chance.
The tight-gap probe, for example, agrees on \HpdFractightgapfiveAcc\ of pairs.
A prompt-held-out fusion of the probes reaches \FusHpdProbespreregisteredAcc\ by assigning the opposite signs.
HPSv2.1 agrees with those annotators at \HpdHpstwoAcc\ on data it was trained on.
ImageReward agrees with ImageRewardDB annotators at \IrImagerewardAcc, with a wide interval and only \NIrTsprompts\ tattoo-subject prompts.

The same probes, read one at a time, detect the injected defects they were defined for, and the reward models detect many of those defects too, without a physical unit.
A Pick-a-Pic slice of \PapHpstwoN\ pairs does not overturn the HPDv2 result, and PickScore on it is in-distribution.
A Gaussian scale-and-spread stress test, with no additive pixel noise, keeps skeleton $F_1$ at \ssFoneSizeThreeSpreadZero\ when the spread is $\sigma{=}0$ and lowers it as $\sigma$ grows.
It remains an assumption, not a skin study.

For AI tattoo generation, a preference reward is evidence about annotators.
A millimetre probe is evidence about a declared needle, gap and print size.
Reporting those probes separately, and stating when a reward model is scored in distribution, is the audit.
